\section{Direction changes and improved occupation}
\label{ang:section}

We now improve the occupation estimate when both particles have high energy
and their relative angle is intermediate. The additional information is a
bound on the number of direction changes of a prescribed angular size.
Its proof uses the signed bootstrap and the baseline estimates already
established; it does not use the improved occupation estimate proved at
the end of this section.

We retain the notation of the preceding sections. In particular,
$L=\log(2+P)$, $S(w)=P^2L/w$, and $m=\min(p,w)$. All constants below may
depend on the datum, the fixed time horizon, and the fixed decomposition
parameters. A constant written $C(M,A)$ may also depend on the two bootstrap
parameters. The lower threshold for $P$ may depend on $M,A$.
Fix a bootstrap endpoint $\tau$: throughout this section the bootstrap is
assumed only on $[0,\tau]$, and all time intervals are clipped to this range.

\Needspace{18\baselineskip}
\subsection{A bound on changes of direction}

\begin{lemma}[Direction changes]\label{ang:direction}
Assume the signed increment bootstrap \eqref{eq:bootstrap} on $[0,\tau]$.
Let $w$ be a dyad with $w\ge P^{3/5}$, and let $J\subset[0,\tau]$ be a closed
interval of length at most $(w/P)^2$ on which a supported characteristic
satisfies $w/8\le q_X\le8w$. Suppose
\[
 w^{-3/5}\le\delta\le w^{-1/10}.
\]
Starting at the left endpoint of $J$, stop whenever the direction
$\xi=V_X/|V_X|$ first has chordal distance $\delta$ from its direction at
the preceding stopping time. Include the last, possibly incomplete,
interval. For sufficiently large $P$, the number of intervals is at most
\begin{equation}\label{ang:count}
 C(M,A)L^6/\delta.
\end{equation}
\end{lemma}

\begin{proof}
Put $I=(w/P)^2$. Throughout $J$, $|V_X|\asymp w$, so normalization of
nonzero vectors gives
\begin{equation}\label{ang:normalization}
 |\xi(t)-\xi(s)|\le Cw^{-1}|V_X(t)-V_X(s)|.
\end{equation}
Define
\[
 \ell=\min\left\{\left(\frac{w\delta}{MP}\right)^2,
                         \frac{w\delta}{AS(w)}\right\}.
\]
The bootstrap and \eqref{ang:normalization} bound the direction change
on an interval of length $c\ell$ by
$C(\sqrt c+c)\delta$. Choosing $c>0$ sufficiently small proves that
each complete first-exit interval has length at least $c\ell$.
In particular their number, denoted by $N_c$, is finite.

On each complete interval choose an endpoint-zero Lipschitz function
$0\le\psi\le1$ with monotone ramps of length $c'\ell$ and value one
between the ramps. Choose $c'$ sufficiently small that the ramps are
disjoint. On a ramp, integrate the removed weight against $\xi_t$ after
subtracting the direction at its outer endpoint. The derivative of the
removed weight has $L^1$ norm one, so the removed signed integral is bounded by
$C(\sqrt{c'}+c')\delta$. Here the bootstrap bounds endpoint-relative
increments; no bound on absolute direction variation is used.
After reducing the fixed constant $c'$ once more, each complete interval
satisfies
\begin{equation}\label{ang:weighted-exit}
 \left|\int\psi\,w\xi_t\,dt\right|\ge\frac12w\delta,
 \qquad \int|\psi_t|\,dt=2.
\end{equation}

Let $P_X=\Id-\xi\otimes\xi$ and
$\Pi=wP_X/|V_X|$, so $w\xi_t=\Pi K$.
Apply Proposition~\ref{can:weighted-formula}
with this multiplier and with the fixed selection
\begin{equation}\label{ang:selection}
 \theta<\kappa\phi,\qquad p\theta\sqrt h\le1.
\end{equation}
The same representative index set is used on every interval. The
weights have disjoint interiors and sum to at most one. Thus, after
taking norms, all bulk errors can be bounded over their union with
the single upper length $I$. Only derivatives of the weights will
produce a factor $N_c$.

First consider the receiver and time-weight derivatives. Their bounds
explain the radial threshold in \eqref{ang:selection}. The signed
primitive contains
$k_0=u-ne/D$, with
\[
 |k_0|\le C\theta/\phi,\qquad |P_Xk_0|\le C\theta.
\]
The projected receiver-acceleration and receiver-cutoff kernels have
size at most
\[ Cf\theta|K|/(wr\phi). \]
The moving projector has precisely
the same size: writing $\rho=|V_X|$ gives
\[
 \Pi_t=-\frac{w}{\rho^2}
 \left[(\xi\cdot K)P_X+(P_XK)\otimes\xi
                         +\xi\otimes(P_XK)\right],
 \qquad
 |\Pi_tk_0|\le\frac{C\theta|K|}{w\phi}.
\]
The representative $w$ is fixed here; no dyad indicator is differentiated.

At fixed receiver time, integrating these kernels over a bin uses
the radial factor $Ch^2$, the velocity density bound $Cp^3\theta^2$,
and an angular cap of area $C\phi^2$. Their coefficient on $|K|$ is
therefore
\begin{equation}\label{ang:receiver-coefficient}
 \frac{Ch^2(p\theta)^3\phi}{w}
 \le \frac{C\sqrt h\,\phi}{w}
 \quad\hbox{on selected indices}.
\end{equation}
The sums in $h,\phi$ are geometric. The remaining $p,\theta$ count
is $O(L^2)$, so the total coefficient is $CL^2/w$.
The baseline absolute-force bound on our interval is
\[
 \int_J|K|\,dt
 \le C\sqrt w\{P\sqrt I+(M+\sqrt A)IS(w)\}
 \le C(M,A)w^{3/2}L.
\]
Thus all receiver and projector derivative terms contribute at most
$C(M,A)w^{1/2}L^3$. They do not contribute a factor $N_c$.

For a weight derivative, the kernel is $f\Pi k_0/r$, and its fixed-time
coefficient is instead
\begin{equation}\label{ang:weight-coefficient}
 Ch^2(p\theta)^3\phi^2\le C\sqrt h\,\phi^2.
\end{equation}
Its sum is $CL^2$. Equation \eqref{ang:weighted-exit} consequently
gives a total weight-derivative cost $CL^2N_c$. Initial-cone boundary
terms are bounded by $C\int\sum\psi\,dt\le CI$ and create no further
cell count.

We now estimate the remaining terms. On the complementary small-sector
bins, $p\theta\sqrt h>1$ will make the projected good-field far bound
summable. Outside the small-scale deficit exception, the baseline
occupation estimates give
\begin{equation}\label{ang:baseline-R}
 \mathsf M\le ChR^2/p,
 \qquad R=1+\sqrt{I/\lambda}.
\end{equation}
For $h>I$ this follows instead from the elementary energy bound
$\mathsf M\le Ch/p$. We use the unrefined $R$ in
\eqref{ang:baseline-R}; in particular no enhanced occupation estimate
enters this proof. The definition of $\lambda$ in the small sector yields
\begin{align}
 \sqrt I\,Rp\phi
 &\le \frac{wp\phi}{P}
   +C\frac{w^2p}{Pm}\left(M+\sqrt{AL\phi}\right)
 \le C(M,A)w\sqrt L,                                      \label{ang:R-bound}
\end{align}
where $wp\le mP$ was used in the last inequality.

The baseline direct force bound contributes at most $C(M,A)wL$.
For the small-sector good-field term,
perpendicular projection supplies an additional factor $\phi$. Its far
bound outside \eqref{ang:selection} is therefore
\[
 C\sqrt I\,R\frac{\phi}{\theta\sqrt h}
 \le \frac{C(M,A)w\sqrt L}{p\theta\sqrt h}.
\]
When the second condition in \eqref{ang:selection} fails, the denominator
is greater than one and $h>P^{-2}$. Hence only $O(L)$ radial dyads occur.
There are also $O(L)$ momentum dyads and $O(L)$ dyads for each angle.
Summing gives at most $C(M,A)wL^5$.

The small deficit exception is entirely within
\eqref{ang:selection} after increasing the fixed constant $C_1$ in the
definition of $\lambda$. Indeed, when $h\le\lambda$ and
$\phi\le C_0/m$,
\[
 p\theta\sqrt h\le p\theta\sqrt\lambda
 \le C\kappa(p/P)m\phi^2/C_1
 \le C\kappa C_0^2/(C_1m)\le1.
\]
Source cutoff derivatives are assigned to overlapping unselected bins
by Lemma~\ref{can:cutoff-cancellation}. At the structural edge
$\theta\asymp\kappa\phi$, use the direct bounds for the remaining sector.
Fixed-factor neighboring indices change only constants in these
estimates.

By \eqref{can:spatial-cutoff-bound}, the projected spatial cutoff kernel is bounded by
$Cf\phi^2/(r^2\theta)$. Its baseline near/far balance, followed by the
geometric sum in $h$, is
\[
 C\sqrt I\,Rp\phi^2
\]
per remaining indices. By \eqref{ang:R-bound}, the sum over $p,\theta,
\phi$ is at most $C(M,A)wL^4$. The projected central residual is
$Cf\phi/r^2$, which is smaller in this sector. For the exceptional
deficit scales, use their near bounds directly. Summing
$h\le\lambda\le C/P^2$ and then $\theta\le C\phi$ gives
\[
 CI P^{-2}p^3\phi^5\le CI p^3/(P^2m^5).
\]
These terms fit the same bulk bound after summing the remaining dyads.
This separate near estimate avoids using \eqref{ang:baseline-R} on its
exceptional set.

Summing \eqref{ang:weighted-exit} now yields the conservative bound
\[
 \tfrac12N_cw\delta
 \le C(M,A)wL^6+CL^2N_c.
\]
Since $w\delta\ge w^{2/5}\ge P^{6/25}$, the last term is absorbed
for sufficiently large $P$. Division by $w\delta$ and addition of
the last incomplete interval prove \eqref{ang:count}.
\end{proof}

\subsection{Occupation at intermediate angles}

The baseline cell argument keeps each velocity within a small multiple
of $\phi$ of its value at a reference hit, so its cell length $\lambda$
shrinks with that tolerance. We now use cells whose length is chosen to
keep only the energies comparable to $p$ and $w$. Within each such cell,
Lemma~\ref{ang:direction} partitions the source and receiver motions into
intervals of nearly fixed direction. Their common refinement permits
the same monotone-crossing argument as before; the number of direction
intervals controls the total hit duration.

\begin{lemma}[Direction-based occupation at intermediate angles]\label{ang:improved-hit}
Assume the signed increment bootstrap on $[0,\tau]$.
Let a small-sector bin satisfy
\begin{equation}\label{ang:window}
 p,w\ge P^{7/10},\qquad
 j^{-49/100}\le\phi\le j^{-4/25}\quad(j=p,w).
\end{equation}
Let the eligible receiver times lie in a closed interval $J\subset[0,\tau]$ of length at most
$I\le(w/P)^2$, with receiver energy in $[w/8,8w]$ at those times.
Set $I_m^0=(m/P)^2$ and $\beta=29/25$. Then, for
$P\ge P_*(M,A)$, the hit measure satisfies
\begin{equation}\label{ang:improved-occupation}
 \mathsf M\le C\frac hp\phi^{-\beta}
      \max\left(1,\frac I{I_m^0},\frac h{I_m^0}\right).
\end{equation}
The constant $C$ is independent of $M,A$.
\end{lemma}

\begin{proof}
The window \eqref{ang:window} implies
\begin{equation}\label{ang:margins}
 m\phi\ge m^{51/100},\qquad
 m^{-2}/\phi\le m^{-151/100},\qquad
 \phi\le P^{-14/125}.
\end{equation}
Thus the bins are actual for large $P$, and speed-magnitude deficits
are small compared to $\phi$.

Choose a sufficiently large fixed constant $C_2$ and put
\[
 Q_0=C_2(M^2+AL),\qquad \lambda_0=I_m^0/Q_0.
\]
For $h>\lambda_0$, the elementary energy estimate already suffices.
Indeed, for sufficiently large $P$, \eqref{ang:margins} implies
$\phi^{-\beta}\ge Q_0$, and hence $h\phi^{-\beta}\ge I_m^0$.
Consequently
\[
 h\phi^{-\beta}\max(1,I/I_m^0,h/I_m^0)\ge\max(I,h),
\]
which proves \eqref{ang:improved-occupation} in this case.

Suppose now that $h\le\lambda_0$. For upper bounds we may enlarge to all
eligible receiver times in $J$ and use closed bin support ranges.
Partition source time into intervals
of length $\lambda_0$, clipped to $[0,\tau]$. All hit
source times lie in an interval of length at most $C(I+h)$, so at most
$C(1+(I+h)/\lambda_0)$ original intervals can contribute.
Fix one such interval $F$. Extend it to the right by $C_T h$, where
$C_T$ bounds $r/h$, and clip this extension to $[0,\tau]$.

If $F$ has a hit, the receiver energy is comparable to $w$ throughout
the extended interval. For every label hitting $F$, its source energy
is comparable to $p$ throughout $F$. To see this, apply the bootstrap
up to the first exit from a fixed enlarged range around a reference
hit energy. On a time interval of length $C_T\lambda_0$, the two
bootstrap terms are at most $Cm/\sqrt{C_2}$ and $Cm/C_2$.
Increasing $C_2$ excludes the first exit. The argument works in both
time directions and is unaffected by clipping.

For each use of Lemma~\ref{ang:direction}, choose a dyad $j'$ nearest
the relevant reference-hit energy. It is comparable to $w$ for the
receiver and to $p$ for a source. The entire interval stays within
$[j'/8,8j']$. Increasing $C_2$ also ensures that its length is at most
$(j'/P)^2$. For large $P$, $j'\ge P^{3/5}$. Choose a small fixed
$c_2>0$ and use $\delta=c_2\phi$. Uniformly for $j'\asymp j$,
\[
 \frac{\delta}{j'^{-3/5}}\ge c j^{11/100}\longrightarrow\infty,
 \qquad
 \frac{\delta}{j'^{-1/10}}\le Cj^{-3/50}\longrightarrow0.
\]
Therefore all the hypotheses of Lemma~\ref{ang:direction} hold.

Choose one common receiver direction partition on the extended
interval. Its number of cells is at most $C(M,A)L^6/\phi$.
For each hit label, independently choose its source direction
partition on $F$, with the same bound on the number of cells.
We shall prove that the hit duration for each such label is at most
\begin{equation}\label{ang:label-duration}
 C(M,A)hL^6/\phi.
\end{equation}

Call a hit bad if a boundary $b$ of the receiver direction partition
lies in $[s,t]$, and call all other hits good.
Since $0<t-s\le C_T h$, the source time of a bad hit belongs
to $[b-C_T h,b]$. The total duration of bad hits is thus bounded by
$C(M,A)hL^6/\phi$. This bound uses only the time difference $t-s$;
it is independent of the retarded derivative $dt/ds$.

At a remaining hit, $s,t$ lie in one receiver direction cell.
Throughout $[s,t]$ its direction varies by at most $2c_2\phi$,
and its speed differs from one by $O(w^{-2})$. The bin condition
$|a(t)-n|\le C\phi$ and \eqref{ang:margins} therefore give
\begin{equation}\label{ang:small-separation}
 |Y(s)-X(s)|
 =\left|\int_s^t(a(\zeta)-n)\,d\zeta\right|
 \le Ch\phi.
\end{equation}

On $F$, refine the source and receiver direction partitions together.
The number of resulting intervals is bounded by the sum of the two
numbers of cells plus one. In one such interval containing a good
hit, choose that hit as reference. Actual small-angle geometry gives
$|u(s)-a(t)|\ge c\phi$. Replacing $a(t)$ by $a(s)$ changes this vector
by at most $Cc_2\phi+Cm^{-2}$. At any other time in the combined
interval, the change of relative velocity has the same bound, since
each direction varies by at most $2c_2\phi$ and both speed deficits
are $O(m^{-2})$. Choose $c_2$ small and then $P$ large. There is a
fixed unit vector $e_F$ for this label and combined interval such that
\[
 e_F\cdot(u(s)-a(s))\ge c\phi
\]
throughout it. Thus $e_F\cdot(Y-X)$ is strictly increasing with
derivative at least $c\phi$. At good hits it lies in an interval of
length at most $Ch\phi$, by \eqref{ang:small-separation}. Their total
duration in this combined interval is consequently at most $Ch$.
This remains true when the good-hit set has several components.
Summing the combined intervals and adding the bad-hit bound proves
\eqref{ang:label-duration}.

The set of labels hitting $F$ has mass at most $C/p$: at any one time
of $F$, all of these labels have energy at least $cp$, by the preceding
stability argument, and total particle energy is bounded. The closed
support convention makes the set of hit labels a compact projection
before the null collision labels are discarded,
as in the baseline occupation argument. The hit set is measurable,
and the preceding reasoning bounds each of its
source-time sections. The label-dependent reference hits and direction
partitions serve only to prove this numerical bound; no measurability
of those choices is required.

Multiplying the mass bound by \eqref{ang:label-duration} and then
counting the original intervals yields
\begin{align}
 \mathsf M
 &\le C(M,A)\frac hp\phi^{-1}L^6
                   \left(1+\frac{I+h}{\lambda_0}\right)\notag\\
 &\le C(M,A)\frac hp\phi^{-1}L^7
                   \max\left(1,\frac I{I_m^0},\frac h{I_m^0}\right).
                                                        \label{ang:before-absorption}
\end{align}
Clipped intervals and final direction remainders satisfy all the same
estimates. Finally, since $\beta-1=4/25$, \eqref{ang:margins} gives
\[
 \phi^{-(\beta-1)}\ge P^{56/3125}.
\]
This positive power absorbs the factor $C(M,A)L^7$ for
$P\ge P_*(M,A)$. Equation \eqref{ang:before-absorption} becomes
\eqref{ang:improved-occupation}, with a constant independent of the
bootstrap parameters.
\end{proof}

The gain in Lemma~\ref{ang:improved-hit} comes from controlling how often
relative motion can change its direction. Within each combined direction
interval, a hit is confined to a short monotone crossing. Only after
this estimate has been proved will it be used in the final selection of
bins for the signed momentum bound.
